\documentclass{beamer}
\usetheme{epyt}
\usepackage{macros_beamer}
\usepackage[portuguese]{babel}
\title{Introdução à Teoria Espectral de Grafos}
\subtitle{Primeiros resultados sobre a matriz de adjacência}
\date{22 de maio de 2026}
\author{Leonardo Suarez}
\institute{Instituto de Matemática e Estatística - UFRGS}

\usepackage[backend=biber, style=ieee]{biblatex}\addbibresource{referencias_grafos.bib}
\setbeamertemplate{bibliography item}{\insertbiblabel}

\DeclareMathOperator{\tr}{tr}
\DeclareMathOperator{\spect}{spect}

\begin{document}

\nocite{notas_col,cvetkovic}

\frame{\titlepage}

\begin{frame}
    \tableofcontents
\end{frame}

\section{O espectro de um grafo}

\begin{frame}{Definição do espectro de um grafo}
    Seja \(G\) um grafo com \(n\) vértices e \(A(G)\) sua matriz de adjacência. O \alert{polinômio característico} de \(G\) é o polinômio \(p_G(\lambda) = \det(\lambda I - A(G))\). O \alert{espectro} de \(G\) é a lista 
    \[\spect(G) = \begin{bmatrix}
        \lambda_1 & \lambda_2 & \cdots & \lambda_r \\
        m_1 & m_2 & \cdots & m_r
        \end{bmatrix},
    \] onde \(\lambda_1, \ldots, \lambda_r\) são os autovalores distintos de \(A(G)\) e \(m_1, m_2, \ldots, m_r\) são suas respectivas multiplicidades. O maior autovalor \(\lambda_1\) é chamado de \alert{índice} de \(G\).
    
\end{frame}

\begin{frame}{Autovetor como valoração de vértices}
    Os \alert{autovalores} de \(A\) são os números reais que satisfazem \(Ax = \lambda x\) para algum vetor \(x \neq 0\). Tal vetor \(x\) é chamado de \alert{autovetor} de \(A\) associado ao autovalor \(\lambda\). A relação \(Ax = \lambda x\) pode ser entendida da seguinte forma: se \(x = (x_1, x_2, \ldots, x_n)^T\) então

    \begin{equation}
        \lambda x_u = \sum_{v \sim u} x_v, \quad u = 1, 2, \ldots, n,
    \end{equation} onde a soma é feita sobre os vizinhos \(v\) de \(u\).
\end{frame}

\begin{frame}
    \begin{itemize}
        \item Sendo \(A(G)\) uma matriz real simétrica, todos seus autovalores são reais. Como \(\tr(A(G)) = 0\), a soma dos autovalores de \(A(G)\) é nula.
        \item Se \(\lambda\) é um autovalor de \(A\) então \(\mathcal{E}(\lambda) = \{x \in \mathbb{R}^n; Ax = \lambda x\}\) é um subespaço vetorial de \(\mathbb{R}^n\) chamado de \alert{autoespaço} de \(G\) associado a \(\lambda\).
        \item O Teorema Espectral assegura que existe uma base ortonormal de \(\mathbb{R}^n\) formada por autovetores de \(A\). Então os autoespaços de \(G\) são dois a dois ortogonais. Além disso, a dimensão de \(\mathcal{E}(\lambda)\) é igual à multiplicidade de \(\lambda\) como raiz de \(p_G(\lambda)\). Um autovalor \emph{simples} é aquele cuja multiplicidade é \(1\).
    \end{itemize}
\end{frame}

\section{Alguns resultados}

\begin{frame}
    \begin{prop}
        Se um grafo \(G\) tem grau máximo \(\Delta(G)\) então todo autovalor \(\lambda\) de \(G\) satisfaz \(\vert \lambda \vert \leq \Delta(G)\).
    \end{prop}
    \begin{proof}[Esboço da prova]
        Toma vértice \(u\) tal que \(\vert x_u \vert\) é máximo e usa a desigualdade triangular.
    \end{proof}
    \pause 
    \begin{prop}
        O número de cadeias de comprimento \(l\) ligando os vértices \(v_i\) e \(v_j\) em um grafo \(G\) é dado por \((A^l)_{ij}\).
    \end{prop}
    \begin{proof}[Esboço da prova]
        Prova por indução em \(l\).
    \end{proof}
\end{frame}

\section{O espectro de certos tipos de grafos}

\begin{frame}
    \begin{prop} \label{prop:polinomio_uniao}
        Se \(G\) é a união de dois grafos \(G_1\) e \(G_2\) então \(p_G(\lambda) = p_{G_1}(\lambda) p_{G_2}(\lambda)\). Em particular, se \(G_1, G_2, \ldots, G_r\) são as componentes conexas de \(G\) então \(\displaystyle p_G(\lambda) = \prod_{i=1}^r p_{G_i}(\lambda)\).
    \end{prop}
    \begin{proof}[Esboço da prova]
        Escreve \(A(G)\) em blocos e calcula o determinante usando a expansão de Laplace.
    \end{proof}
\end{frame}

\begin{frame}[allowframebreaks]
    \begin{prop}
        Seja \(G\) um grafo regular de grau \(k\). Então
        \begin{enumerate}
            \item [a.] \(k\) é um autovalor de \(G\);
            \item [b.] \(G\) é conexo se e somente se \(k\) é um autovalor simples de \(G\);
            \item [c.] qualquer autovalor \(\lambda\) de \(G\) satisfaz \(|\lambda| \leq k\).
        \end{enumerate}
    \end{prop}
    \begin{proof}[Esboço da prova]
        \begin{enumerate}
            \item [a.] Perceba que a soma das entradas da linha \(j\) da matriz de adjacência fornece o grau do vértice \(v_j\).
            \item [b.] \((\Rightarrow)\) Seja \(x = [x_1 \; \ldots \; x_n]^T\) autovetor de \(G\) associado ao autovalor \(k\). Prova que \(x\) é múltiplo escalar do vetor \(1\).
    
            \((\Leftarrow)\) Suponha que \(G\) é desconexo, com componentes conexas \(k\)-regulares \(G_1, \ldots, G_m\). Usa a proposição anterior para mostrar que \(k\) é autovalor de \(G\) de multiplicidade \(>1\).
            \item [c.] Aplica a mesma ideia do item anterior com \(y = [y_1 \; \ldots \; y_n]^T\) autovetor de \(G\) associado ao autovalor \(\lambda\).
        \end{enumerate}
    \end{proof}
\end{frame}

\begin{frame}[allowframebreaks]
    \begin{prop}
        Seja \(G\) um grafo \(k\)-regular com \(n\) vértices e autovalores \(\lambda_1 = k, \lambda_2, \ldots, \lambda_n\). Então os autovalores de \(\overline{G}\) são \(n-1-k, -1-\lambda_2, \ldots, -1-\lambda_n\), respectivamente, associados aos mesmos autovetores de \(G\).
    \end{prop}
    \begin{proof}[Esboço da prova]
        \begin{itemize}
            \item Verifica que \(A(\overline{G}) = J - I - A(G)\), onde \(J\) é a matriz de \(n \times n\) cujas entradas são todas iguais a \(1\).
            \item Sendo \(\{1, u_2, \ldots, u_n\}\) uma base ortonormal de autovetores de \(A(G)\) associada aos autovalores \(k, \lambda_2, \ldots, \lambda_n\), respectivamente, aplica \(A(\overline{G})\) a cada um desses vetores.
            \item Mostra que \(1\) é autovetor de \(A(\overline{G})\) associado ao autovalor \(n-1-k\) e que \(u_i\) é autovetor de \(A(\overline{G})\) associado ao autovalor \(-1-\lambda_i\), para \(i = 2, \ldots, n\).
        \end{itemize}
    \end{proof}
\end{frame}

\begin{frame}
    \begin{exmp}
        Se \(G\) é um grafo sem arestas com \(n\) vértices, então \(G\) é \(0\)-regular com \(p_G(\lambda) = \lambda^n\) e \[\spect(G) = \begin{bmatrix}
            0 \\
            n
        \end{bmatrix}.\]
    \end{exmp}
    \pause
    \begin{exmp}
        O grafo completo com \(n\) vértices \(K_n\) é o complementar do grafo sem arestas. Temos então \(p_{K_n}(\lambda) = (\lambda - (n-1))(\lambda + 1)^{n-1}\) e \[\spect(K_n) = \begin{bmatrix}
            n-1 & -1 \\
            1 & n-1
        \end{bmatrix}.\]
    \end{exmp}
\end{frame}

\begin{frame}
    \begin{exmp}
        Se \(G\) é o grafo formado pela união de \(s\) cópias de \(K_2\) então \(G\) é \(1\)-regular e tem \(2s\) vértices, com \(p_{K_2} (\lambda) = (\lambda - 1)(\lambda + 1)\), vale \(p_G (\lambda) = (\lambda - 1)^s(\lambda + 1)^s\) e \[ \spect(G) = \begin{bmatrix}
            1 & -1 \\
            s & s
        \end{bmatrix}.\] \pause 
        O grafo complementar de \(G\) é chamado de \emph{hiperoctaedro} (ou \emph{cocktail party graph}) e denotado por \(H_s\). O espectro de \(H_s\) é \[\spect(H_s) = \begin{bmatrix}
            2s-1 & -1 \\
            1 & 2s-1
        \end{bmatrix}.\]
    \end{exmp}
\end{frame}

\begin{frame}
    \begin{exmp}
    Seja \(G\) um grafo bipartido cujo conjunto de vértices é particionado em \(V_1\) e \(V_2\) com \(|V_1| = k_1\) e \(|V_2| = k_2\). Se rotulamos primeiro os vértices de \(V_1\) e depois os de \(V_2\), a matriz de adjacência de \(G\) tem a forma
    \[
        A(G) = \begin{bmatrix}
            0 & B \\
            B^T & 0
        \end{bmatrix},
    \] onde \(B\) é uma matriz de adjacência de dimensões \(k_1 \times k_2\). Afirmamos que \(\lambda\) é um autovalor de \(G\) se e somente se \(-\lambda\) é um autovalor de \(G\).
    \end{exmp}
\end{frame}

\section{Isomorfismo e coespectralidade de grafos}

\begin{frame}{Grafo sem arestas é determinado pelo seu espectro}
    \begin{prop}
        Um grafo \(G\) possui um único autovalor se e somente se \(G\) é totalmente desconexo, ou seja, é um grafo sem arestas.
    \end{prop}
    \begin{proof}[Esboço da prova]
        Já provamos a direção \((\Leftarrow)\). Para \((\Rightarrow)\), mostra que \(\tr(A(G)) = 0 \implies \lambda = 0\) e que \(A(G)\) é a matriz nula.
    \end{proof}
\end{frame}

\begin{frame}
    \begin{prop}
        Se \(G\) tem exatamente dois autovalores distintos \(\lambda_1 > \lambda_2\), com multiplicidades \(m_1\) e \(m_2\), respectivamente, então \(G\) é o grafo \(\lambda_1\)-regular formado por \(m_1\) cópias de \(K_{\lambda_1 + 1}\). Além disso, temos \(\lambda_2 = -1\) e \(m_2 = m_1\lambda_1\).
    \end{prop}

    \begin{proof}[Esboço da prova]
        \begin{itemize}
            \item Avalia \(m(A) = (A - \lambda_1 I)(A - \lambda_2 I)\) e considera os elementos \((A^2)_{kk}\) (número de arestas partindo de \(v_k\)).
            \item Mostra que \(G\) é \(\lambda_1\)-regular e que \(\lambda_2 = -1\).
            \item Mostra que \(G\) é a união de \(m_1\) cópias de \(K_{\lambda_1 + 1}\) e que \(m_2 = m_1\lambda_1\).
        \end{itemize}
    \end{proof}
\end{frame}

\section*{Referências}

\begin{frame}
    \printbibliography
\end{frame}

\end{document}